% Figure 1 -- Newton polygon of an irregular singularity (p_0, p_1, p_2)
\begin{tikzpicture}[x=1pt,y=1pt,scale=1.12,transform shape,>=stealth,line width=0.5pt]
  \def\u{14.145}
  % axes
  \draw[->] (0,0) -- (0,187.3);
  \draw[->] (0,0) -- (187.3,0);
  \node[inner sep=0pt] at (-2.2,202.5) {\fontsize{15}{18}\selectfont$z$};
  \node[inner sep=0pt] at (209.1,2.7)  {\fontsize{15}{18}\selectfont$x$};
  % lattice points on the boundary
  \foreach \i in {1,...,7} {\fill (0,{\i*\u}) circle[radius=2.7pt];}
  \foreach \i in {1,...,9} {\fill ({\i*\u},0) circle[radius=2.7pt];}
  % Newton polygon
  \draw (0,{7*\u}) -- ({3*\u},{6*\u}) -- ({6*\u},{4*\u}) -- ({9*\u},0);
  % labels of the boundary points
  \node[inner sep=0pt,font=\fontsize{10.6}{13}\selectfont] at (51.4,99.0)  {$p_2$};
  \node[inner sep=0pt,font=\fontsize{10.6}{13}\selectfont] at (106.3,62.7) {$p_1$};
  \node[inner sep=0pt,font=\fontsize{10.6}{13}\selectfont] at (138.1,16.8) {$p_0$};
\end{tikzpicture}
